\documentclass[11pt]{exam}
\usepackage{tikz}
\usetikzlibrary {plotmarks}
\usepackage{multicol}
\usepackage{subfigure}

\newcommand{\class}{MATH 96}
\newcommand{\term}{Spring 2025}
\newcommand{\examnum}{Exam 1} 
\newcommand{\examdate}{2/19/25}
\newcommand{\timelimit}{60 Minutes}

\parindent 0ex

\begin{document} 

\pagestyle{head}
\firstpageheader{}{}{}
\runningheader{\class}{\examnum\ - Page \thepage\ of \numpages}{\examdate}
\runningheadrule

\begin{flushright}
\begin{tabular}{p{2.2in} r l}
\textbf{\class} & \textbf{Name:}     & \makebox[2.7in]{\hrulefill}\\
\textbf{\term}  &                    &\\
\textbf{\examnum} &  & \\
 & \textbf{Section:} & \makebox[2.7in]{\hrulefill}\\

\end{tabular}\\
\end{flushright}
\rule[1ex]{\textwidth}{.1pt}

This exam contains \numpages\ pages (including this cover page) and
\numquestions\ problems.  Check to see if any pages are missing.  Enter
your name and your section above.
\vspace{0.2in}

You may \emph{not} use textbooks, course notes, cell phones, or calculators during this exam.
\vspace{0.2in}

You are required to show your work on each problem on this exam, and the following rules apply:
\vspace{0.2in}

\begin{minipage}[t]{3.7in}
\vspace{0pt}
\begin{itemize}

\item \textbf{Organize your work in a reasonable, neat, and coherent way.} Work that is jumbled, disorganized, or lacks clear reasoning will receive little credit.  

\item \textbf{Partial credit may be given.} Any answer must be supported by calculations, explanation, and/or algebraic work to receive full credit.  Unsupported answers will not receive full credit, but well-argued, incorrect solutions may earn partial credit.

\item \textbf{If you require more space, use the back of a page.} Clearly indicate when you do this below the problem.
\end{itemize}

Do not write in the table to the right.
\end{minipage}
\hfill
\begin{minipage}[t]{2.3in}
\vspace{0pt}
\gradetablestretch{2}
\vqword{Problem}
\addpoints
\gradetable[v][pages]
\end{minipage}
%=======================================================================================================================
\newpage
\begin{questions}
%=======================================================================================================================
\question[15] Simplify and reduce each expression to a single fraction.
\begin{multicols}{2}
\begin{parts}
\noaddpoints
\part[8] $\displaystyle \left(\frac{2}{3}-\frac{1}{4}\right) \div\left(\frac{3}{4}+\frac{1}{3}\right)$
\vspace{2.25in}
\part[7] $\displaystyle  \frac{-3-(-1^2)}{-(-2)^2+(3-2)^2}$
\end{parts}
\addpoints
\end{multicols}
\vspace{3.5in}
\question[15] Solve the equation $\displaystyle \frac{5}{4}+\frac{4}{3}x = \frac{1}{4}$
\vspace{4in}
\addpoints

\question[10] Simplify the expression as much as possible.
\[
- (x+2y) - 2(-3y-2x) + 2y
\]
\vspace{3.5in}

\question[15] Find the equation of the line in slope-intercept form that passes through the points $(1,-1)$ and $(-1,3)$.
\vspace{3.5in}

\newpage
\question[20] Consider the line with the equation $-2x + 4y = -8.$
\noaddpoints
\begin{parts}
\part[5] Find the slope of the line above.
\vspace{1in}
\part[5] Find the $x$-intercept and $y$-intercept of the line above.
\vspace{1in}
\part[5] Determine whether the line $\displaystyle y - 1 = -2(x+2)$ is parallel, perpendicular, or neither to the line above.  Explain your reasoning.
\vspace{1.75in}
\part[5] Graph the line $-2x + 4y =-8$. \emph{Clearly label any two points with $P$ and $Q$ on the line.}
\begin{figure}[ht]
\vspace{-12pt}
\centering
\begin{tikzpicture}[scale=.9]
\draw[dotted] (-7.5,-4.5) grid (7.5,4.5); 
\draw[<->] (-7.5,0) -- (7.5,0)node[anchor=south] {\footnotesize{$x$}};
\draw[<->] (0,-4.5)--(0,4.5)node[anchor=west] {\footnotesize{$y$}};
\foreach \x in {-7,-6,-5,-4,-3,-2,-1,1,2,3,4,5,6,7}
    \draw (\x cm, 2pt) -- (\x cm, -2pt) node[anchor=north] {\footnotesize{$\x$}};
\foreach \y in {-4,-3,-2,-1,1,2,3,4}
    \draw (2pt, \y cm) -- (-2pt, \y cm) node[anchor=east] {\footnotesize{$\y$}}; 
\end{tikzpicture}
\end{figure}
\end{parts}
\addpoints
\newpage
%=======================================================================================================================

\vspace{2.75in}
\emph{\textbf{Clearly define your variables. Solve each problem by setting up an algebraic equation and then use that equation to find the solution. There’s no need to simplify any numerical fractions.}}
\question[10] Find three consecutive even integers whose sum is 222.

\vspace{2.75in}
\question[15] A store is offering a 25\% discount on a jacket that originally costs \$240. 

\begin{parts}
    
\noaddpoints
    \part[10] What is the amount of the discount?
    \vspace{3in}
    \part[5] What is the sale price of the jacket after the discount is applied?
\end{parts}
\addpoints
%\newpage
%==============================================================================================================
%\emph{Clearly identify your variables.  Solve each problem by setting up an equation algebraically.  Then use the equation to solve the problem.  You do not need to simplify any numeric fraction.}
%\question[12] Tony Stark flew on an airplane from New York City to London. It took him 5
%hours to fly with constant wind, and 7 hours on the return flight against the same constant
%wind. If the speed of the wind is 40 mph, find the speed of the plane in still air.
%=======================================================================================================================
\end{questions}
\end{document}


}
